% =====================================================================
\chapter{Soluciones de los ejercicios}
\label{ap:soluciones}
% =====================================================================

Cada solución se presenta en cinco partes: datos, planteamiento, desarrollo, resultado y comprobación. En los pasos intermedios se redondea a cuatro cifras significativas salvo que se indique otra cosa; los recuentos de parámetros y de bytes son enteros exactos. Se usan las unidades binarias de la presentación: $1$~KiB $=1024$~bytes, $1$~MiB $=1024^{2}$~bytes y $1$~GiB $=1024^{3}=1\,073\,741\,824$~bytes.

% ---------------------------------------------------------------------
\section{Tema 1}

\begin{solucion}{eje:1-vf}
\datos tres afirmaciones sobre RMSNorm, SwiGLU y la regla de las cabezas.
\planteamiento contrastar cada una con la \cref{prop:rmsnorm}, la \cref{prop:presupuesto} y la \cref{tab:hiper}.
\pasos
\begin{enumerate}[label=(\alph*)]
  \item \textbf{Falsa.} RMSNorm no resta la media. Con $x=(1,0)$ y $t=1$: $\RMS(x)=\sqrt{1/2}$ y $\RMSNorm(x)=(1{,}414;\ 0)$, mientras que $x+t\uno=(2,1)$ tiene $\RMS=\sqrt{5/2}=1{,}581$ y $\RMSNorm(x+t\uno)=(1{,}265;\ 0{,}6325)$. La invariancia ante traslaciones es propia de LayerNorm.
  \item \textbf{Verdadera.} SwiGLU con $\dff=4d$ tiene $3\cdot d\cdot 4d=12d^{2}$ parámetros, frente a $2\cdot d\cdot 4d=8d^{2}$: $12/8=1{,}5$. Por eso se reduce $\dff$ a $\frac83 d$.
  \item \textbf{Verdadera.} Llama~3 8B: $32\cdot 128=4096$; 70B: $64\cdot 128=8192$; 405B: $128\cdot 128=16\,384$; Mixtral: $32\cdot 128=4096$; Grok-1: $48\cdot 128=6144$.
\end{enumerate}
\resultado \textbf{(a) falsa; (b) verdadera; (c) verdadera.}
\comprobacion en (a), $\norm{(1{,}265;\ 0{,}6325)}^{2}=1{,}600+0{,}4000=2=d$, como exige la \cref{prop:rmsnorm}(i).
\end{solucion}

\begin{solucion}{eje:1-norm}
\datos $x=(2,-1,0,3)$, $d=4$, $\gamma=\uno$, $\beta=0$, $\varepsilon=0$.
\planteamiento aplicar la \cref{def:rmsnorm}; para $3x$, la \cref{prop:rmsnorm}(ii).
\pasos
\begin{enumerate}[label=\arabic*.]
  \item $\sum_i x_i^{2}=4+1+0+9=14$; $\RMS(x)=\sqrt{14/4}=\sqrt{3{,}5}=1{,}871$.
  \item $\RMSNorm(x)=(2,-1,0,3)/1{,}871=(1{,}069;\ -0{,}5345;\ 0;\ 1{,}604)$.
  \item LayerNorm: $\mu(x)=(2-1+0+3)/4=1$; $c=x-\uno=(1,-2,-1,2)$; $\sum_i c_i^{2}=1+4+1+4=10$; $\RMS(c)=\sqrt{2{,}5}=1{,}581$; $\LN(x)=(0{,}6325;\ -1{,}265;\ -0{,}6325;\ 1{,}265)$.
  \item $\RMS(3x)=3\cdot 1{,}871=5{,}612$, luego $\RMSNorm(3x)=(6,-3,0,9)/5{,}612$, que vale $(1{,}069;\ -0{,}5345;\ 0;\ 1{,}604)=\RMSNorm(x)$.
  \item Norma: $1{,}069^{2}+0{,}5345^{2}+0+1{,}604^{2}=1{,}143+0{,}2857+2{,}571=4{,}000=d$.
\end{enumerate}
\resultado \textbf{$\RMSNorm(x)=\RMSNorm(3x)=(1{,}069;\ -0{,}5345;\ 0;\ 1{,}604)$ y $\LN(x)=(0{,}6325;\ -1{,}265;\ -0{,}6325;\ 1{,}265)$.}
\comprobacion la salida de LayerNorm tiene media $0$ y norma al cuadrado $0{,}4+1{,}6+0{,}4+1{,}6=4=d$, como la de RMSNorm aplicada al vector centrado (\cref{prop:rmsnorm}(iv)).
\end{solucion}

\begin{solucion}{eje:1-llama3}
\datos \cref{tab:hiper}: $(L,d,\dff,H,\nkv)=(32,4096,14\,336,32,8)$, $(80,8192,28\,672,64,8)$ y $(126,16\,384,53\,248,128,8)$.
\planteamiento $\dhd=d/H$; cabezas por grupo $H/\nkv$; bloque $=2d^{2}$ ($W_Q$, $W_O$) $+\,2d\,\nkv\dhd$ ($W_K$, $W_V$) $+\,3d\,\dff$ (SwiGLU) $+\,2d$ (normas).
\pasos
\begin{enumerate}[label=\arabic*.]
  \item Dimensión por cabeza: $4096/32=8192/64=16\,384/128=128$ en los tres.
  \item Consultas por cabeza K/V: $32/8=4$, $64/8=8$ y $128/8=16$.
  \item 8B: atención $2\cdot 4096^{2}+2\cdot 4096\cdot 1024=33\,554\,432+8\,388\,608=41\,943\,040$; MLP $3\cdot 4096\cdot 14\,336=176\,160\,768$; normas $8192$; bloque $218\,112\,000$.
  \item 70B: atención $2\cdot 8192^{2}+2\cdot 8192\cdot 1024=150\,994\,944$; MLP $3\cdot 8192\cdot 28\,672=704\,643\,072$; normas $16\,384$; bloque $855\,654\,400$.
  \item 405B: atención $2\cdot 16\,384^{2}+2\cdot 16\,384\cdot 1024=570\,425\,344$; MLP $3\cdot 16\,384\cdot 53\,248=2\,617\,245\,696$; normas $32\,768$; bloque $3\,187\,703\,808$.
  \item Fracción del MLP en 8B: $176\,160\,768/218\,112\,000=0{,}8077$.
\end{enumerate}
\resultado \textbf{$\dhd=128$ en los tres; 4, 8 y 16 consultas por cabeza K/V; bloques de $218\,112\,000$, $855\,654\,400$ y $3\,187\,703\,808$ parámetros; el MLP es el $80{,}77\,\%$ del bloque de 8B.}
\comprobacion los bloques del modelo 8B suman $32\cdot 218\,112\,000=6{,}980\cdot 10^{9}$; con \emph{embeddings} de entrada y salida de $128\,256\times 4096$ (unos $1{,}05\cdot 10^{9}$) el total queda cerca de los $8\cdot 10^{9}$ del nombre.
\end{solucion}

\begin{solucion}{eje:1-swiglu}
\datos $d=512$; SwiGLU con $\dff\approx\frac83d$ redondeado a múltiplo de~64; SiLU con $\dff=4d$.
\planteamiento SwiGLU tiene $3d\,\dff$ parámetros y el MLP SiLU, $2d\cdot 4d$ (\cref{prop:presupuesto}).
\pasos
\begin{enumerate}[label=\arabic*.]
  \item $\frac83\cdot 512=1365{,}3$; $1365{,}3/64=21{,}33$, cuyo entero más próximo es $21$; $\dff=21\cdot 64=1344$.
  \item SwiGLU: $3\cdot 512\cdot 1344=2\,064\,384$.
  \item SiLU: $2\cdot 512\cdot 2048=2\,097\,152$.
  \item Diferencia relativa: $(2\,097\,152-2\,064\,384)/2\,097\,152=32\,768/2\,097\,152=0{,}015\,625=1{,}563\,\%$.
\end{enumerate}
\resultado \textbf{$\dff=1344$; $2\,064\,384$ frente a $2\,097\,152$ parámetros; el SwiGLU tiene un $1{,}563\,\%$ menos.} La comparación es justa porque ambos MLP tienen casi los mismos parámetros y FLOPs; una ventaja del SwiGLU no puede atribuirse a un mayor tamaño.
\comprobacion sin redondeo, $3\cdot 512\cdot 1365{,}3=2\,097\,152$, exactamente igual al MLP SiLU; toda la diferencia procede del redondeo $1365{,}3\to 1344$.
\end{solucion}

\begin{solucion}{eje:1-inv}
\datos $x\in\R^{d}$ no constante, $t\in\R$, $\varepsilon=0$, $\beta=0$.
\planteamiento usar $\mu(x+t\uno)=\mu(x)+t$ y la \cref{prop:rmsnorm}(iv).
\pasos
\begin{enumerate}[label=\arabic*.]
  \item $\mu(x+t\uno)=\frac1d\sum_i(x_i+t)=\mu(x)+t$, luego $(x+t\uno)-\mu(x+t\uno)\uno=x-\mu(x)\uno$.
  \item Por la \cref{prop:rmsnorm}(iv), $\LN(x+t\uno)=\RMSNorm(x-\mu(x)\uno)=\LN(x)$; $x+t\uno$ tampoco es constante.
  \item Contraejemplo para RMSNorm: $x=(1,-1)$, $t=1$. $\RMSNorm(x)=(1,-1)$, porque $\RMS(x)=1$. $x+t\uno=(2,0)$, $\RMS=\sqrt2$ y $\RMSNorm(x+t\uno)=(1{,}414;\ 0)\ne(1,-1)$.
  \item Recuperación salvo escala: si $\gamma=\uno$, $u=\RMSNorm(x)=x/\RMS(x)$, luego $x=\RMS(x)\,u$ con $\RMS(x)>0$: $x$ es un múltiplo positivo de $u$. Recíprocamente, todo múltiplo positivo de $u$ da la misma salida por la \cref{prop:rmsnorm}(ii).
\end{enumerate}
\resultado \textbf{LayerNorm es invariante ante traslaciones y RMSNorm no; la salida de RMSNorm determina la entrada exactamente salvo un factor positivo.} $\blacksquare$
\comprobacion LayerNorm pierde más información: determina $x$ salvo un factor positivo \emph{y} una traslación, es decir, pierde dos grados de libertad frente a uno de RMSNorm.
\end{solucion}

\begin{solucion}{eje:1-aspecto}
\datos (a) $\dff=\rho d$, $\nkv$ cabezas K/V, $\dhd=d/H$; (b) $N_{\mathrm{b}}=7\cdot 10^{10}$, $r=102$; (c) Llama~3 70B: $L=80$, $d=8192$, $H=64$, $\nkv=8$, $\rho=3{,}5$.
\planteamiento repetir el recuento de la \cref{prop:aspecto} con las nuevas matrices.
\pasos
\begin{enumerate}[label=\arabic*.]
  \item (a) Atención: $W_Q,W_O$ suman $2d^{2}$; $W_K,W_V$ tienen $d\times\nkv\dhd=d\times(\nkv/H)d$ cada una, en total $2(\nkv/H)d^{2}$. MLP: $3\rho d^{2}$. Bloque: $c\,d^{2}$ con $c=2+2\nkv/H+3\rho$. Con $L=d/r$: $N_{\mathrm{b}}=c\,d^{3}/r$ y $d=(rN_{\mathrm{b}}/c)^{1/3}$. Con $\nkv=H$ y $\rho=8/3$, $c=2+2+8=12$.
  \item (b) $d=(102\cdot 7\cdot 10^{10}/12)^{1/3}=(5{,}950\cdot 10^{11})^{1/3}=8411$; $L=8411/102=82{,}46$. Llama~3 70B tiene $d=8192$ y $L=80$.
  \item (c) $c=2+2\cdot 8/64+3\cdot 3{,}5=2+0{,}25+10{,}5=12{,}75$; $N_{\mathrm{b}}=80\cdot 8192^{2}\cdot 12{,}75=68\,451\,041\,280$.
  \item Con $c=12{,}75$, la fórmula de (a) da $d=(102\cdot 7\cdot 10^{10}/12{,}75)^{1/3}=8243$, más cerca de $8192$.
\end{enumerate}
\resultado \textbf{(a) $d=(rN_{\mathrm{b}}/c)^{1/3}$ con $c=2+2\nkv/H+3\rho$; (b) $d\approx 8411$, $L\approx 82$; (c) $c=12{,}75$ y $N_{\mathrm{b}}=68\,451\,041\,280$.} La diferencia de (b) se debe a que Llama~3 ahorra en $W_K,W_V$ ($0{,}25d^{2}$ en lugar de $2d^{2}$) y gasta más en el MLP ($10{,}5d^{2}$ en lugar de $8d^{2}$).
\comprobacion el recuento exacto del \cref{eje:1-llama3}, $80\cdot 855\,654\,400=68\,452\,352\,000$, difiere de (c) en $1\,310\,720=80\cdot 2\cdot 8192$, que son las normalizaciones despreciadas.
\end{solucion}

% ---------------------------------------------------------------------
\section{Tema 2}

\begin{solucion}{eje:2-vf}
\datos tres afirmaciones sobre GQA, MQA y MLA.
\planteamiento contrastar con la \cref{prop:gqa}, la \cref{prop:cache} y la \cref{def:mla}.
\pasos
\begin{enumerate}[label=(\alph*)]
  \item \textbf{Falsa.} GQA conserva las $H$ cabezas de consulta; $W_Q$ sigue siendo $d\times d$ y sus FLOPs no cambian (\cref{prop:gqa}(ii)).
  \item \textbf{Verdadera.} Con $\nkv=1$, la caché por token y capa es $2\dhd s$ bytes, sin $H$ en la fórmula; en bf16 con $\dhd=128$, $512$~bytes.
  \item \textbf{Falsa.} MLA guarda un único vector latente $c_t$ de dimensión $\dc$ por token y capa (más la clave de RoPE desacoplada, compartida), no claves y valores por cabeza.
\end{enumerate}
\resultado \textbf{(a) falsa; (b) verdadera; (c) falsa.}
\comprobacion en (c), con los datos de DeepSeek-V2, MLA guarda $576$ números por token y capa, frente a $2\cdot 128\cdot 128=32\,768$ si guardase claves y valores de todas las cabezas.
\end{solucion}

\begin{solucion}{eje:2-kv}
\datos $L=32$, $H=32$, $\dhd=128$, $s=2$, $T=8192$, $b=1$; $\nkv=32$, $8$ y $1$.
\planteamiento \cref{prop:cache}: $M_{\mathrm{KV}}=2L\nkv\dhd Tbs$.
\pasos
\begin{enumerate}[label=\arabic*.]
  \item MHA por token: $2\cdot 32\cdot 32\cdot 128\cdot 2=524\,288$~bytes; total $524\,288\cdot 8192=4\,294\,967\,296$~bytes $=4$~GiB.
  \item GQA-8 por token: $2\cdot 32\cdot 8\cdot 128\cdot 2=131\,072$~bytes; total $1\,073\,741\,824$~bytes $=1$~GiB.
  \item MQA por token: $2\cdot 32\cdot 1\cdot 128\cdot 2=16\,384$~bytes; total $134\,217\,728$~bytes $=128$~MiB.
\end{enumerate}
\resultado \textbf{MHA 4~GiB; GQA-8 1~GiB; MQA 128~MiB.}
\comprobacion los cocientes son $4/1=4=32/8$ y $4096/128=32=32/1$, es decir, $H/G$ (\cref{prop:gqa}(i)).
\end{solucion}

\begin{solucion}{eje:2-cs336}
\datos $V=32\,000$, $d=5120$, $\dff=13\,824$, $L=40$, $H=40$, $\dhd=128$, $T=1024$, $b=64$, $s=2$, ancho de banda $3{,}35\cdot 10^{12}$~bytes/s.
\planteamiento parámetros $=2Vd+3d\dff L+(2dH\dhd+2d\nkv\dhd)L$; memoria $=2\cdot$parámetros $+\,b\,M_{\mathrm{KV}}$; latencia $=$ memoria$/$ancho de banda; rendimiento $=b/$latencia.
\pasos
\begin{enumerate}[label=\arabic*.]
  \item Partes comunes: $2Vd=327\,680\,000$; $3d\dff L=8\,493\,465\,600$; $2dH\dhd L=2\,097\,152\,000$.
  \item (a) $\nkv=40$: $2d\nkv\dhd L=2\,097\,152\,000$; parámetros $13\,015\,449\,600$; pesos en bf16, $26\,030\,899\,200$~bytes.
  \item Caché por secuencia: $2\cdot 40\cdot 40\cdot 128\cdot 1024\cdot 2=838\,860\,800$~bytes; con $b=64$, $53\,687\,091\,200$~bytes.
  \item Memoria total: $79\,717\,990\,400$~bytes $\approx 79{,}72$~GB; latencia $79{,}72\cdot 10^{9}/3{,}35\cdot 10^{12}=0{,}023\,80$~s $=23{,}80$~ms; rendimiento $64/0{,}023\,80=2689$ tokens/s.
  \item (b) $\nkv=8$: $2d\nkv\dhd L=419\,430\,400$; parámetros $11\,337\,728\,000$; pesos $22\,675\,456\,000$~bytes.
  \item Caché por secuencia: $167\,772\,160$~bytes; con $b=64$, $10\,737\,418\,240$~bytes.
  \item Memoria total $33\,412\,874\,240$~bytes $\approx 33{,}41$~GB; latencia $9{,}974$~ms; rendimiento $6417$ tokens/s.
\end{enumerate}
\resultado \textbf{(a) $1{,}302\cdot 10^{10}$ parámetros, $79{,}72$~GB, $23{,}80$~ms por paso y $2689$ tokens/s; (b) $1{,}134\cdot 10^{10}$ parámetros, $33{,}41$~GB, $9{,}974$~ms y $6417$ tokens/s.} GQA reduce la caché a la quinta parte y, en este régimen limitado por la memoria, multiplica el rendimiento por $2{,}386$.
\comprobacion la caché se divide exactamente por $40/8=5$: $53\,687\,091\,200/10\,737\,418\,240=5$. La clase de CS336 añade que el ahorro permite subir el lote a 256 ($65{,}63$~GB y $13\,068$ tokens/s), mientras que con MHA no cabría en 80~GB.
\end{solucion}

\begin{solucion}{eje:2-intens}
\datos Llama~2 70B: $H=64$; $\nkv=64$ (MHA), $8$ (GQA) o $1$ (MQA); $s=2$; $\Icrit=156$~FLOP/byte.
\planteamiento \cref{prop:intensidad}: $I=2H/(s\nkv)$; según el modelo \emph{roofline}, la fracción del máximo es $\min(1,I/\Icrit)$.
\pasos
\begin{enumerate}[label=\arabic*.]
  \item MHA: $I=2\cdot 64/(2\cdot 64)=1$; fracción $1/156=0{,}6410\,\%$.
  \item GQA-8: $I=128/16=8$; fracción $8/156=5{,}128\,\%$.
  \item MQA: $I=128/2=64$; fracción $64/156=41{,}03\,\%$.
\end{enumerate}
\resultado \textbf{$I=1$, $8$ y $64$~FLOP/byte; como máximo un $0{,}6410\,\%$, un $5{,}128\,\%$ y un $41{,}03\,\%$ del rendimiento de la A100.} Las tres variantes siguen limitadas por la memoria.
\comprobacion los cocientes entre intensidades son $8=H/G$ y $64=H/1$, como afirma la \cref{prop:gqa}(iii).
\end{solucion}

\begin{solucion}{eje:2-mla}
\datos $L=40$, $H=32$, $\dhd=128$, $\dc=256$, $d_h^{R}=64$, $s=2$, presupuesto $16$~GiB.
\planteamiento \cref{tab:mla}: MLA $(\dc+d_h^{R})L$ y MHA $2H\dhd L$ números por token; grupos equivalentes $G=(\dc+d_h^{R})/(2\dhd)$.
\pasos
\begin{enumerate}[label=\arabic*.]
  \item MLA: $(256+64)\cdot 40=12\,800$ números; $25\,600$~bytes $=25$~KiB por token.
  \item MHA: $2\cdot 32\cdot 128\cdot 40=327\,680$ números; $655\,360$~bytes $=640$~KiB por token. Cociente $327\,680/12\,800=25{,}6$.
  \item Grupos equivalentes: $320/256=1{,}25$.
  \item Tokens en $16$~GiB: MLA, $17\,179\,869\,184/25\,600=671\,088{,}64$, es decir, $671\,088$; MHA, $17\,179\,869\,184/655\,360=26\,214{,}4$, es decir, $26\,214$.
\end{enumerate}
\resultado \textbf{(a) 25~KiB frente a 640~KiB por token ($25{,}6$ veces menos); (b) $1{,}25$ grupos; (c) unos $671\,088$ tokens con MLA y $26\,214$ con MHA.}
\comprobacion $671\,088{,}64/26\,214{,}4=25{,}6$, el mismo cociente que en el paso~2.
\end{solucion}

\begin{solucion}{eje:2-pool}
\datos grupo $\mathcal G_g$ de $m=H/G$ cabezas con $W_K^{(i)}=W_K^{\ast}$ y $W_V^{(i)}=W_V^{\ast}$ para $i\in\mathcal G_g$.
\planteamiento la conversión de Ainslie \emph{et al.} asigna al grupo la media de sus proyecciones~\cite[§2.1]{gqa}.
\pasos
\begin{enumerate}[label=\arabic*.]
  \item (a) $W_K^{(g)}=\frac1m\sum_{i\in\mathcal G_g}W_K^{(i)}=\frac1m\,m\,W_K^{\ast}=W_K^{\ast}$, y lo mismo para $W_V$.
  \item Cada cabeza $i$ del grupo calcula las mismas claves $XW_K^{\ast}$ y valores $XW_V^{\ast}$ antes y después de la conversión, y su consulta $XW_Q^{(i)}$ no cambia; las puntuaciones, los pesos y la salida de la cabeza son idénticos.
  \item (b) Sean $d=\dhd=1$, dos cabezas con $W_Q^{(1)}=W_Q^{(2)}=1$, $W_K^{(1)}=1$ y $W_K^{(2)}=3$, una consulta procedente de una entrada igual a~1 (luego $q=1$ en ambas cabezas) y dos posiciones previas con entradas $x_1=1$ y $x_2=2$. Antes de convertir, la cabeza~1 tiene claves $1$ y $2$ y puntuaciones $1$ y $2$; la cabeza~2, claves $3$ y $6$ y puntuaciones $3$ y $6$. Tras promediar, $W_K=(1+3)/2=2$, las claves son $2$ y $4$ en ambas cabezas y las puntuaciones, $2$ y $4$.
\end{enumerate}
\resultado \textbf{(a) si las proyecciones de cada grupo coinciden, la conversión es exacta; (b) en general cambia las puntuaciones, por lo que hace falta el ajuste posterior (\emph{uptraining}).} $\blacksquare$
\comprobacion en (b), las diferencias de puntuación entre las dos posiciones pasan de $1$ y $3$ a $2$ y $2$: los pesos de la softmax cambian en ambas cabezas.
\end{solucion}

\begin{solucion}{eje:2-rope}
\datos vectores fila; $k_{j,i}=c_jW^{UK}_iR_j$, consulta $q_{t,i}R_t$; $R_mR_n\T=R_{m-n}$; frecuencias $0<\theta_l<2\pi$.
\planteamiento desarrollar la traspuesta del producto y analizar cuándo $R_\delta(W^{UK}_i)\T$ no depende de $\delta$.
\pasos
\begin{enumerate}[label=\arabic*.]
  \item (a) $q_{t,i}R_t(c_jW^{UK}_iR_j)\T=q_{t,i}R_tR_j\T(W^{UK}_i)\T c_j\T=q_{t,i}R_{t-j}(W^{UK}_i)\T c_j\T$. Luego $M_{t,j}=R_{t-j}(W^{UK}_i)\T\in\R^{\dhd\times\dc}$.
  \item (b) Si $M_{t,j}$ es la misma para todos los pares, en particular $R_{1}(W^{UK}_i)\T=R_0(W^{UK}_i)\T=(W^{UK}_i)\T$. Cada columna $w$ de $(W^{UK}_i)\T$ cumple $R_1w=w$.
  \item $R_1$ es diagonal por bloques con rotaciones de ángulo $\theta_l$. Una rotación plana de ángulo $\theta_l\notin 2\pi\mathbb Z$ solo fija el vector nulo, porque $R(\theta)w=w$ con $w\ne 0$ exige que $1$ sea autovalor de $R(\theta)$, cuyos autovalores son $e^{\pm i\theta}\ne 1$.
  \item Por tanto, cada par de componentes de $w$ es nulo, $w=0$, y $W^{UK}_i=0$.
  \item Con $W^{UK}_i\ne 0$, la matriz que relaciona $q_{t,i}$ con $c_j$ cambia con la distancia $t-j$; no puede precalcularse una única $W^{Q}_i(W^{UK}_i)\T$ que sirva para todas las posiciones de la caché.
\end{enumerate}
\resultado \textbf{$M_{t,j}=R_{t-j}(W^{UK}_i)\T$, que solo es constante si $W^{UK}_i=0$; por eso MLA añade una clave rotada aparte.} $\blacksquare$
\comprobacion sin RoPE, $R_m=I$ para todo $m$ y $M_{t,j}=(W^{UK}_i)\T$, que es la matriz absorbida en la \cref{prop:absorcion}.
\end{solucion}

% ---------------------------------------------------------------------
\section{Tema 3}

\begin{solucion}{eje:3-vf}
\datos tres afirmaciones sobre la ventana deslizante y la interpolación de posiciones.
\planteamiento contrastar con la \cref{def:swa}, la \cref{prop:circular} y la \cref{def:pi}.
\pasos
\begin{enumerate}[label=(\alph*)]
  \item \textbf{Verdadera.} La caché circular tiene $W$ ranuras por capa (\cref{prop:circular}); su tamaño, $2\nkv\dhd Ws$ bytes, no contiene $T$.
  \item \textbf{Verdadera.} PI divide por $s_{\mathrm{e}}$ el ángulo $m\theta_i$ de todos los pares por igual; el cambio de base NTK divide $\theta_i$ por $s_{\mathrm{e}}^{2i/(\dhd-2)}$, un factor que va de $1$ en el par más rápido a $s_{\mathrm{e}}$ en el más lento (\cref{prop:ntk}).
  \item \textbf{Falsa.} La posición $i$ atiende a $j\in[i-W+1,i]$; la posición $i-W$ queda fuera. Solo puede influir a través de una segunda capa.
\end{enumerate}
\resultado \textbf{(a) verdadera; (b) verdadera; (c) falsa.}
\comprobacion en (c), el \cref{teo:campo} con $L=1$ da un campo receptivo de $W$ posiciones, de $i-W+1$ a $i$.
\end{solucion}

\begin{solucion}{eje:3-pares}
\datos $T=16\,384$, $W=4096$.
\planteamiento \cref{prop:pares}.
\pasos
\begin{enumerate}[label=\arabic*.]
  \item Completa: $T(T+1)/2=16\,384\cdot 16\,385/2=268\,451\,840/2=134\,225\,920$.
  \item Ventana, filas $1$ a $W$: $W(W+1)/2=4096\cdot 4097/2=8\,390\,656$.
  \item Filas $W+1$ a $T$: $(T-W)W=12\,288\cdot 4096=50\,331\,648$.
  \item Total con ventana: $8\,390\,656+50\,331\,648=58\,722\,304$.
  \item Cociente: $134\,225\,920/58\,722\,304=2{,}286$; la aproximación da $T/(2W)=2$.
\end{enumerate}
\resultado \textbf{$134\,225\,920$ frente a $58\,722\,304$ pares; cociente $2{,}286$, frente a $2$ de la aproximación.}
\comprobacion con la fórmula cerrada, $TW-W(W-1)/2=67\,108\,864-8\,386\,560=58\,722\,304$. La aproximación $T/(2W)$ subestima porque $T/W=4$ no es grande; con $T=32\,768$ el cociente exacto es $4{,}267$ y la aproximación, $4$.
\end{solucion}

\begin{solucion}{eje:3-circular}
\datos $W=4$; posiciones $0,\dots,9$.
\planteamiento \cref{def:circular}: la posición $i$ va a la ranura $i\bmod 4$.
\pasos
\begin{enumerate}[label=\arabic*.]
  \item Ranuras: $0\to 0$, $1\to 1$, $2\to 2$, $3\to 3$, $4\to 0$, $5\to 1$, $6\to 2$, $7\to 3$, $8\to 0$, $9\to 1$.
  \item Contenido final: la ranura 0 guarda la última posición escrita en ella, la 8; la ranura 1, la 9; la ranura 2, la 6; la ranura 3, la 7.
  \item La posición 9 necesita $j\in[9-4+1,9]=\{6,7,8,9\}$, que es exactamente el contenido de la caché.
\end{enumerate}
\resultado \textbf{tras la posición 9, las ranuras $(0,1,2,3)$ contienen las posiciones $(8,9,6,7)$, que son las de su ventana.}
\comprobacion las posiciones $6,7,8,9$ tienen restos $2,3,0,1$ módulo 4, todos distintos, como exige la demostración de la \cref{prop:circular}.
\end{solucion}

\begin{solucion}{eje:3-ntk}
\datos $\dhd=128$, $\beta_0=10^{4}$, $s_{\mathrm{e}}=4$.
\planteamiento \cref{def:ntk} y \cref{prop:ntk}: $\beta_0'=\beta_0s_{\mathrm{e}}^{128/126}$ y $\theta_i'=\theta_is_{\mathrm{e}}^{-2i/126}$; PI divide todas por $s_{\mathrm{e}}$.
\pasos
\begin{enumerate}[label=\arabic*.]
  \item $s_{\mathrm{e}}^{128/126}=4^{1{,}016}=4{,}089$; $\beta_0'=40\,890$.
  \item Originales: $\theta_0=1$; $\theta_{32}=10^{4\cdot(-64/128)}=10^{-2}=0{,}01$; $\theta_{63}=10^{-4\cdot 126/128}=1{,}155\cdot 10^{-4}$.
  \item PI: $0{,}25$; $0{,}0025$; $2{,}887\cdot 10^{-5}$.
  \item NTK: $\theta_0'=1$; $\theta_{32}'=0{,}01\cdot 4^{-64/126}=0{,}01/2{,}022=0{,}004\,945$; $\theta_{63}'=1{,}155\cdot 10^{-4}/4=2{,}887\cdot 10^{-5}$.
  \item (c) Para $\theta_{32}$, PI divide por $4$ y NTK por $2{,}022\approx\sqrt{4}$.
  \item (d) $\sqrt{1/t}=0{,}1\ln 4+1=0{,}1\cdot 1{,}386+1=1{,}139$.
\end{enumerate}
\resultado \textbf{$\beta_0'=40\,890$; con NTK, $\theta_0$ no cambia, $\theta_{63}$ se divide por 4 como en PI y $\theta_{32}$ se divide por $2{,}022$ en lugar de por 4; $\sqrt{1/t}=1{,}139$.}
\comprobacion el par central ($i=32$, a mitad del recorrido $0\dots 63$) se divide por $s_{\mathrm{e}}^{64/126}\approx s_{\mathrm{e}}^{1/2}$: el cambio de base reparte la interpolación geométricamente entre $1$ y $s_{\mathrm{e}}$.
\end{solucion}

\begin{solucion}{eje:3-d2l}
\datos $n=4096$, $d=1024$, $k=3$.
\planteamiento sustituir en $n^{2}d$, $nd^{2}$ y $knd^{2}$; el cociente atención$/$recurrente es $n/d$.
\pasos
\begin{enumerate}[label=\arabic*.]
  \item Atención: $4096^{2}\cdot 1024=17\,179\,869\,184\approx 1{,}718\cdot 10^{10}$.
  \item Recurrente: $4096\cdot 1024^{2}=4\,294\,967\,296\approx 4{,}295\cdot 10^{9}$.
  \item Convolución: $3\cdot 4\,294\,967\,296=12\,884\,901\,888\approx 1{,}288\cdot 10^{10}$.
  \item $n^{2}d>nd^{2}\iff n>d=1024$.
  \item La convolución conecta posiciones a distancia $k-1$ por capa; para relacionar dos posiciones a distancia $n$ hacen falta del orden de $n/k$ capas, mientras que la atención las conecta en una. La red recurrente necesita $O(n)$ pasos secuenciales.
\end{enumerate}
\resultado \textbf{$1{,}718\cdot 10^{10}$, $4{,}295\cdot 10^{9}$ y $1{,}288\cdot 10^{10}$; la atención es más cara que la recurrente para $n>1024$.} La ventaja de la atención es su camino de longitud uno entre cualquier par de posiciones y su paralelismo; su inconveniente, el coste cuadrático que motiva la ventana deslizante y los SSM.
\comprobacion el cociente atención$/$recurrente es $1{,}718\cdot 10^{10}/4{,}295\cdot 10^{9}=4=n/d$.
\end{solucion}

\begin{solucion}{eje:3-dilatada}
\datos ventana dilatada: la posición $i$ lee $i-a\delta$ con $a\in\{0,\dots,w-1\}$; $L$ capas.
\planteamiento inducción sobre $L$, describiendo el conjunto de desplazamientos alcanzables $D_L=\{i-j\}$.
\pasos
\begin{enumerate}[label=\arabic*.]
  \item $D_1=\{0,\delta,\dots,(w-1)\delta\}$; el máximo es $\delta(w-1)$.
  \item Si $D_L$ es el conjunto de desplazamientos tras $L$ capas, la capa $L+1$ suma un desplazamiento de $D_1$: $D_{L+1}=\{u+v:u\in D_1,\ v\in D_L\}$.
  \item Por inducción, $D_L=\{\delta(a_1+\dots+a_L):a_l\in\{0,\dots,w-1\}\}$, cuyo máximo es $L\delta(w-1)$, alcanzado con todos los $a_l=w-1$. La salida en $i$ solo depende de $[i-L\delta(w-1),i]$ y el extremo sí influye.
  \item Todos los elementos de $D_L$ son múltiplos de $\delta$. Si $\delta>1$, las posiciones $i-1,\dots,i-\delta+1$ no influyen: el intervalo no se cubre entero.
\end{enumerate}
\resultado \textbf{el campo receptivo abarca $L\delta(w-1)+1$ posiciones, pero con $\delta>1$ solo influyen las que distan un múltiplo de $\delta$.} La estimación $L\delta w$ de Longformer es la longitud del intervalo; para cubrir los huecos, el artículo usa dilataciones distintas en distintas cabezas. $\blacksquare$
\comprobacion con $\delta=1$ se recupera el \cref{teo:campo}: $L(w-1)+1$ posiciones, todas alcanzables.
\end{solucion}

\begin{solucion}{eje:3-ntk-dem}
\datos $\theta_i'=(\beta_0')^{-2i/\dhd}$; PI usa el ángulo $(m/s_{\mathrm{e}})\theta_i$ con $0\le m<L'=s_{\mathrm{e}}L_0$.
\planteamiento (a) imponer la condición sobre $\theta'_{\dhd/2-1}$ y despejar $\beta_0'$; (b) acotar el ángulo.
\pasos
\begin{enumerate}[label=\arabic*.]
  \item (a) $\theta_0'=(\beta_0')^{0}=1=\theta_0$ para cualquier base: la primera condición se cumple siempre.
  \item $\theta'_{\dhd/2-1}=(\beta_0')^{-(\dhd-2)/\dhd}$ y se exige que valga $\beta_0^{-(\dhd-2)/\dhd}/s_{\mathrm{e}}$.
  \item Elevando ambos miembros a $-\dhd/(\dhd-2)$ (la función $u\mapsto u^{-\dhd/(\dhd-2)}$ es biyectiva en $(0,\infty)$): $\beta_0'=\beta_0\,s_{\mathrm{e}}^{\dhd/(\dhd-2)}$, única solución.
  \item (b) Con PI, el mayor ángulo del par $i$ se alcanza en $m=L'-1$: $(L'-1)\theta_i/s_{\mathrm{e}}=(L_0-1/s_{\mathrm{e}})\theta_i<L_0\theta_i$.
  \item En el entrenamiento, el mayor ángulo es $(L_0-1)\theta_i$; con PI, $(L_0-1/s_{\mathrm{e}})\theta_i$, que solo excede al anterior en $(1-1/s_{\mathrm{e}})\theta_i<\theta_i$, menos de un paso de posición.
\end{enumerate}
\resultado \textbf{(a) $\beta_0'=\beta_0s_{\mathrm{e}}^{\dhd/(\dhd-2)}$ es la única base con esas dos propiedades; (b) con PI los ángulos quedan dentro del intervalo $[0,L_0\theta_i)$, es decir, a menos de un paso del máximo visto en el entrenamiento.} $\blacksquare$
\comprobacion con $\dhd=128$ y $s_{\mathrm{e}}=4$, (a) reproduce $\beta_0'=40\,890$ del \cref{eje:3-ntk}.
\end{solucion}

% ---------------------------------------------------------------------
\section{Tema 4}

\begin{solucion}{eje:4-memoria}
\datos (a) $N=8$, $k=2$; (b) un token con una sola asignación y experto sin capacidad.
\planteamiento \cref{prop:activos} y \cref{def:capacidad}.
\pasos
\begin{enumerate}[label=\arabic*.]
  \item (a) El enrutador puede elegir cualquier experto en cualquier token y capa, de modo que los 8 deben estar en memoria; cada token usa 2.
  \item Parámetros de MLP guardados$/$usados: $8P_E/(2P_E)=4$. El modelo denso con los mismos activos guarda solo $2P_E$ por capa.
  \item (b) El experto lleno no evalúa ese token; la contribución del MLP es nula y la salida del bloque es $h+0=h$.
  \item El vector del token atraviesa la capa por la conexión residual con toda la información previa; solo se pierde la transformación de ese MLP en esa capa.
\end{enumerate}
\resultado \textbf{(a) la MoE guarda 4 veces más parámetros de MLP de los que usa por token; (b) el token se descarta en esa capa, pero su información sigue en el flujo residual.}
\comprobacion con $k=N$ el cociente es $N/N=1$: sin dispersión no hay sobrecoste de memoria.
\end{solucion}

\begin{solucion}{eje:4-vf}
\datos tres afirmaciones sobre Gemini~1.5, el Switch Transformer y DeepSeek-V3.
\planteamiento contrastar con la \cref{tab:moe} y las observaciones del tema.
\pasos
\begin{enumerate}[label=(\alph*)]
  \item \textbf{Falsa.} El informe solo dice que es un Transformer MoE disperso; no publica el número de expertos~\cite[§3.1]{gemini15}. La cifra de 8 corresponde a Mixtral y Grok-1.
  \item \textbf{Verdadera.} El Switch Transformer usa enrutado top-1~\cite[§2.1]{switch}.
  \item \textbf{Falsa.} DeepSeek-V3 sustituye la pérdida por expertos por sesgos ajustados en línea, pero conserva una pérdida de balanceo por secuencia con un coeficiente muy pequeño~\cite[§2.1.2]{deepseekv3}.
\end{enumerate}
\resultado \textbf{(a) falsa (no es público); (b) verdadera; (c) falsa.}
\comprobacion la clase~4 de CS336 describe el método de DeepSeek-V3 como balanceo «sin pérdida auxiliar» y advierte de que no lo es del todo~\cite{cs336}.
\end{solucion}

\begin{solucion}{eje:4-router}
\datos $z=(0{,}5;\ 2;\ 1{,}2;\ 0)$, $k=2$.
\planteamiento elegir los dos logits mayores y aplicar la softmax entre ellos (\cref{def:moe}); para la comparación, la softmax sobre los cuatro.
\pasos
\begin{enumerate}[label=\arabic*.]
  \item Orden: $2$ (experto 2) $>1{,}2$ (experto 3) $>0{,}5>0$. $\mathcal T=\{2,3\}$.
  \item $e^{2}=7{,}389$; $e^{1{,}2}=3{,}320$; suma $10{,}71$.
  \item $g_2=7{,}389/10{,}71=0{,}6900$; $g_3=3{,}320/10{,}71=0{,}3100$. Equivalentemente, $g_2=1/(1+e^{-0{,}8})=1/1{,}449=0{,}6900$.
  \item Softmax sobre los cuatro: $e^{0{,}5}+e^{2}+e^{1{,}2}+e^{0}=1{,}649+7{,}389+3{,}320+1=13{,}36$; $p_2=0{,}5532$ y $p_3=0{,}2486$.
\end{enumerate}
\resultado \textbf{expertos 2 y 3, con $g_2=0{,}6900$ y $g_3=0{,}3100$; sin renormalizar serían $0{,}5532$ y $0{,}2486$.}
\comprobacion $p_2/(p_2+p_3)=0{,}5532/0{,}8018=0{,}6900=g_2$, como establece la \cref{prop:renorm}.
\end{solucion}

\begin{solucion}{eje:4-bal}
\datos (a) $N=4$, $f=(0{,}4;\ 0{,}1;\ 0{,}35;\ 0{,}15)$, $P=(0{,}35;\ 0{,}15;\ 0{,}35;\ 0{,}15)$; (b) $B_{\mathrm{t}}=2048$, $N=16$, $k=1$, $c=1{,}0$.
\planteamiento \cref{def:bal} y \cref{def:capacidad}.
\pasos
\begin{enumerate}[label=\arabic*.]
  \item Productos: $0{,}4\cdot 0{,}35=0{,}14$; $0{,}1\cdot 0{,}15=0{,}015$; $0{,}35\cdot 0{,}35=0{,}1225$; $0{,}15\cdot 0{,}15=0{,}0225$.
  \item Suma: $0{,}3$; $\Lbal=4\cdot 0{,}3=1{,}2$.
  \item (b) Asignaciones: $kB_{\mathrm{t}}=2048$; reparto ideal $2048/16=128$; $C_{\mathrm{exp}}=1{,}0\cdot 128=128$.
  \item Las 16 capacidades suman $2048$, exactamente las asignaciones. Si un experto recibe más de 128 tokens, otro recibe menos; los sobrantes del primero se descartan mientras quedan plazas libres en el segundo.
\end{enumerate}
\resultado \textbf{(a) $\Lbal=1{,}2$; (b) $C_{\mathrm{exp}}=128$, sin holgura: cualquier desequilibrio provoca descartes.}
\comprobacion en (a), $f$ y $P$ suman 1 y el valor está en $[1,4]$, aunque $f\ne P$; con $f=P$ el \cref{teo:bal} lo garantizaría.
\end{solucion}

\begin{solucion}{eje:4-grok}
\datos $d=6144$, $L=64$, $H=48$, $\nkv=8$, $\dhd=128$, $N=8$, $k=2$, $V=131\,072$, \emph{embeddings} compartidos, $\dff=\lfloor 8d\cdot 2/3\rfloor$, tres matrices por experto.
\planteamiento \cref{prop:activos}: $P_0$ = atención $+$ \emph{embeddings} $+$ enrutadores.
\pasos
\begin{enumerate}[label=\arabic*.]
  \item $\dff=\lfloor 49\,152\cdot 2/3\rfloor=32\,768$ (ya es múltiplo de 8).
  \item $P_E=3\cdot 6144\cdot 32\,768=603\,979\,776$; todos los expertos: $8\cdot 64\cdot P_E=309\,237\,645\,312$.
  \item Atención por capa: $2\cdot 6144^{2}+2\cdot 6144\cdot 1024=75\,497\,472+12\,582\,912=88\,080\,384$; en 64 capas, $5\,637\,144\,576$.
  \item \emph{Embeddings} compartidos: $131\,072\cdot 6144=805\,306\,368$; enrutadores: $64\cdot 6144\cdot 8=3\,145\,728$.
  \item $\Ntot=309\,237\,645\,312+5\,637\,144\,576+805\,306\,368+3\,145\,728=315\,683\,241\,984$.
  \item $\Nact=2\cdot 64\cdot P_E+5\,637\,144\,576+805\,306\,368+3\,145\,728=77\,309\,411\,328+6\,445\,596\,672=83\,755\,008\,000$.
  \item Fracción activa: $83\,755\,008\,000/315\,683\,241\,984=26{,}53\,\%$.
\end{enumerate}
\resultado \textbf{$P_E=603\,979\,776$; $\Ntot\approx 3{,}157\cdot 10^{11}$ y $\Nact\approx 8{,}376\cdot 10^{10}$ ($26{,}53\,\%$).} El total supera en un $0{,}54\,\%$ la cifra redondeada del archivo \texttt{README} ($3{,}14\cdot 10^{11}$); xAI no explica el redondeo, de modo que la diferencia no puede atribuirse con seguridad a ninguna parte del modelo. La fracción activa es coherente con «una cuarta parte de los pesos por token».
\comprobacion los expertos son el $97{,}96\,\%$ de los parámetros,\[309\,237\,645\,312/315\,683\,241\,984=0{,}9796,\]y por eso $\Nact/\Ntot$ queda cerca de $k/N=0{,}25$.
\end{solucion}

\begin{solucion}{eje:4-cota}
\datos $f,P$ distribuciones sobre $\{1,\dots,N\}$.
\planteamiento acotar $\sum_if_iP_i$ por debajo con la no negatividad y por arriba con $P_i\le\max_jP_j\le 1$.
\pasos
\begin{enumerate}[label=\arabic*.]
  \item $f_iP_i\ge 0$, luego $\Lbal\ge 0$; se alcanza con $f=(1,0,\dots,0)$ y $P=(0,1,0,\dots,0)$.
  \item $\sum_if_iP_i\le\max_jP_j\sum_if_i\le 1$, luego $\Lbal\le N$; se alcanza con $f=P=(1,0,\dots,0)$.
  \item En la práctica $f$ y $P$ no son independientes: con top-1, un token se asigna al experto $i$ solo si $p_i(x)$ es la mayor de las $N$ probabilidades, y entonces $p_i(x)\ge 1/N$. Por tanto $P_i\ge\frac1T\sum_{x\to i}p_i(x)\ge f_i/N$.
  \item De ahí, $\Lbal=N\sum_if_iP_i\ge\sum_if_i^{2}\ge 1/N>0$ (Cauchy--Schwarz, como en el \cref{teo:bal}): el valor $0$ no es alcanzable por un enrutador real.
\end{enumerate}
\resultado \textbf{$0\le\Lbal\le N$ para distribuciones arbitrarias, con ambas cotas alcanzables; pero, si $f$ procede del enrutador, $\Lbal\ge 1/N$, y el valor de referencia es el del reparto uniforme, $\Lbal=1$.} $\blacksquare$
\comprobacion con $N=4$ y reparto uniforme, $\sum_if_i^{2}=1/4$ y $\Lbal=1\ge 1/4$.
\end{solucion}

\begin{solucion}{eje:4-grad}
\datos $T$ tokens, top-1, $\Lbal=N\sum_if_iP_i$ con $P_i=\frac1T\sum_xp_i(x)$; (b) $T=8$, $N=4$, $\alpha=10^{-2}$, asignaciones $(4,2,1,1)$.
\planteamiento $f$ es constante; $P_i$ depende linealmente de $p_i(x)$.
\pasos
\begin{enumerate}[label=\arabic*.]
  \item (a) $\partial P_j/\partial p_i(x)=\frac1T$ si $j=i$ y $0$ en otro caso.
  \item $\partial(\alpha\Lbal)/\partial p_i(x)=\alpha Nf_i\cdot\frac1T=\frac{\alpha N}{T}\cdot\frac1T\sum_{x'}\uno[\arg\max_jp_j(x')=i]=\frac{\alpha N}{T^{2}}\sum_{x'}\uno[\cdot]$.
  \item (b) $\alpha N/T^{2}=0{,}04/64=6{,}25\cdot 10^{-4}$; derivadas: $4\cdot 6{,}25\cdot 10^{-4}=2{,}5\cdot 10^{-3}$, $1{,}25\cdot 10^{-3}$, $6{,}25\cdot 10^{-4}$ y $6{,}25\cdot 10^{-4}$.
  \item (c) El descenso de gradiente reduce $p_i(x)$ en proporción a esa derivada: el experto con más asignaciones ve bajar su probabilidad con más fuerza. Como la softmax suma~1, la masa pasa a los demás, y en el lote siguiente el reparto tiende a igualarse.
\end{enumerate}
\resultado \textbf{la derivada es $\alpha N/T^{2}$ por el número de tokens asignados al experto; en el ejemplo, $(2{,}5;\ 1{,}25;\ 0{,}625;\ 0{,}625)\cdot 10^{-3}$.} Es una realimentación negativa: la penalización crece con el tráfico. $\blacksquare$
\comprobacion la derivada del experto 1 es cuatro veces la de los expertos 3 y 4, igual que su número de asignaciones ($4/1$).
\end{solucion}

% ---------------------------------------------------------------------
\section{Tema 5}

\begin{solucion}{eje:5-scan}
\datos S4 con $(\Abar,\Bbar,C)$ fijos; Mamba con $B_t$, $C_t$ y $\Delta_t$ dependientes de la entrada.
\planteamiento la forma convolucional exige un núcleo independiente de $t$ (\cref{prop:conv}); el barrido exige asociatividad (\cref{prop:scan}).
\pasos
\begin{enumerate}[label=\arabic*.]
  \item En S4, $y_t=\sum_jC\Abar^{\,j}\Bbar x_{t-j}$ con un núcleo $\bar K_j=C\Abar^{\,j}\Bbar$ fijo: es una convolución y se calcula en paralelo.
  \item En Mamba, el coeficiente de $x_{t-j}$ es $C_t\Abar_t\cdots\Abar_{t-j+1}\Bbar_{t-j}$, que depende de $t$: no existe un núcleo común.
  \item La recurrencia sigue siendo afín en el estado: cada paso es $h\mapsto\Abar_th+\Bbar_tx_t$. La composición de afines es afín y asociativa (\cref{prop:scan}), y un barrido paralelo calcula todos los prefijos con profundidad $O(\log T)$.
\end{enumerate}
\resultado \textbf{S4 es invariante en el tiempo y admite convolución; Mamba no, pero su linealidad en el estado permite el barrido paralelo.}
\comprobacion si en Mamba se fijan $\Abar_t=\Abar$ y $\Bbar_t=\Bbar$, el coeficiente pasa a ser $C\Abar^{\,j}\Bbar$ y se recupera el núcleo de S4.
\end{solucion}

\begin{solucion}{eje:5-conv}
\datos $\Abar=0{,}9$, $\Bbar=C=1$, $x=(1,0,0,0,1)$, $h_0=0$.
\planteamiento recurrencia $h_t=0{,}9h_{t-1}+x_t$; núcleo $\bar K_j=0{,}9^{j}$.
\pasos
\begin{enumerate}[label=\arabic*.]
  \item $h_1=1$; $h_2=0{,}9$; $h_3=0{,}81$; $h_4=0{,}729$; $h_5=0{,}9\cdot 0{,}729+1=0{,}6561+1=1{,}6561$.
  \item Núcleo: $\bar K=(1;\ 0{,}9;\ 0{,}81;\ 0{,}729;\ 0{,}6561)$.
  \item Convolución en $t=5$: $\bar K_0x_5+\bar K_4x_1=1+0{,}6561=1{,}6561$; para $t\le 4$ solo actúa $x_1$ y $y_t=\bar K_{t-1}$.
\end{enumerate}
\resultado \textbf{$y=(1;\ 0{,}9;\ 0{,}81;\ 0{,}729;\ 1{,}6561)$ por ambos métodos.}
\comprobacion la entrada inicial pesa $0{,}9^{4}=0{,}6561$ cuatro pasos después: la memoria decae más despacio que con $\Abar=0{,}5$ en el \cref{ej:ssm}.
\end{solucion}

\begin{solucion}{eje:5-zoh}
\datos $a=-2$, $\beta=1$, $\Delta=0{,}5$.
\planteamiento \cref{prop:zoh}; bilineal $\Abar=(1+\Delta a/2)/(1-\Delta a/2)$, $\Bbar=\Delta\beta/(1-\Delta a/2)$; Euler $\Abar=1+\Delta a$, $\Bbar=\Delta\beta$.
\pasos
\begin{enumerate}[label=\arabic*.]
  \item Orden cero: $\Abar=e^{-1}=0{,}3679$; $\Bbar=(e^{-1}-1)/(-2)=0{,}6321/2=0{,}3161$.
  \item Bilineal: $\Abar=(1-0{,}5)/(1+0{,}5)=0{,}3333$; $\Bbar=0{,}5/1{,}5=0{,}3333$.
  \item Euler: $\Abar=1-1=0$; $\Bbar=0{,}5$.
  \item Estabilidad para cualquier $\Delta$: con orden cero, $\Abar=e^{a\Delta}\in(0,1)$. Con la bilineal, $\abs{1+\Delta a/2}<\abs{1-\Delta a/2}$ porque $a<0$, luego $\abs{\Abar}<1$. Con Euler, $\abs{1+\Delta a}<1$ exige $\Delta<2/\abs{a}=1$: con $\Delta=1{,}5$, $\Abar=-2$ y la recurrencia diverge.
\end{enumerate}
\resultado \textbf{orden cero $(0{,}3679;\ 0{,}3161)$, bilineal $(0{,}3333;\ 0{,}3333)$, Euler $(0;\ 0{,}5)$; los dos primeros son estables para todo $\Delta>0$ y Euler solo para $\Delta<1$.}
\comprobacion con entrada constante $x=1$, el equilibrio del sistema continuo es $h^{\ast}=-\beta/a=0{,}5$, y los tres métodos lo conservan: $0{,}3161/(1-0{,}3679)=0{,}5$, $0{,}3333/0{,}6667=0{,}5$ y $0{,}5/1=0{,}5$.
\end{solucion}

\begin{solucion}{eje:5-mit1}
\datos (a) $y[n]=0{,}75y[n-1]+x[n-1]$, $x=\delta$, reposo; (b) etapas $g[n]=\frac12g[n-1]+x[n]$ y $h[n]=\frac14h[n-1]+g[n]$.
\planteamiento aplicar las recurrencias desde el reposo.
\pasos
\begin{enumerate}[label=\arabic*.]
  \item (a) $y[0]=0{,}75\cdot 0+x[-1]=0$; $y[1]=0+x[0]=1$; $y[2]=0{,}75$; $y[3]=0{,}5625$; $y[4]=0{,}4219$.
  \item Fórmula: $y[n]=0{,}75^{\,n-1}$ para $n\ge 1$ e $y[0]=0$: un SSM con $\Abar=0{,}75$ y la entrada retrasada un paso.
  \item (b) Primera etapa: $g=(1;\ 0{,}5;\ 0{,}25;\ 0{,}125)$.
  \item Segunda etapa: $h[0]=1$; $h[1]=0{,}25+0{,}5=0{,}75$; $h[2]=0{,}1875+0{,}25=0{,}4375$; $h[3]=0{,}109\,375+0{,}125=0{,}2344$.
  \item Fórmula: $2\cdot 1-1=1$; $2\cdot 0{,}5-0{,}25=0{,}75$; $2\cdot 0{,}25-0{,}0625=0{,}4375$; $2\cdot 0{,}125-0{,}015\,625=0{,}2344$.
\end{enumerate}
\resultado \textbf{(a) $(0;\ 1;\ 0{,}75;\ 0{,}5625;\ 0{,}4219)$, $y[n]=0{,}75^{\,n-1}$; (b) $(1;\ 0{,}75;\ 0{,}4375;\ 0{,}2344)$, que coincide con la fórmula.}
\comprobacion la cascada es la convolución de los núcleos: $h[2]=\sum_{j=0}^{2}(1/2)^{j}(1/4)^{2-j}=1/16+1/8+1/4=7/16=0{,}4375$.
\end{solucion}

\begin{solucion}{eje:5-puerta}
\datos $h_0=1$, $z=(1,-1,3)$, $x=(0,2,-1)$.
\planteamiento \cref{teo:puerta}: $\Delta_t=\softplus(z_t)$, $g_t=\sigma(z_t)$, $h_t=(1-g_t)h_{t-1}+g_tx_t$.
\pasos
\begin{enumerate}[label=\arabic*.]
  \item $t=1$: $\Delta_1=\ln(1+e)=1{,}313$; $g_1=\sigma(1)=0{,}7311$; $h_1=0{,}2689\cdot 1+0{,}7311\cdot 0=0{,}2689$.
  \item Comprobación por orden cero en $t=1$: $\Abar_1=e^{-1{,}313}=0{,}2689$ y $\Bbar_1=1-0{,}2689=0{,}7311$; $h_1=0{,}2689$.
  \item $t=2$: $\Delta_2=\ln(1+e^{-1})=0{,}3133$; $g_2=0{,}2689$; $h_2=0{,}7311\cdot 0{,}2689+0{,}2689\cdot 2=0{,}1966+0{,}5379=0{,}7345$.
  \item $t=3$: $\Delta_3=\ln(1+e^{3})=3{,}049$; $g_3=0{,}9526$; $h_3=0{,}047\,43\cdot 0{,}7345+0{,}9526\cdot(-1)=0{,}034\,83-0{,}9526=-0{,}9177$.
\end{enumerate}
\resultado los pasos son $\Delta=(1{,}313;\ 0{,}3133;\ 3{,}049)$ y las puertas, $g=(0{,}7311;\ 0{,}2689;\ 0{,}9526)$; los estados son \textbf{$h=(0{,}2689;\ 0{,}7345;\ -0{,}9177)$}.
\comprobacion un paso $\Delta_t$ grande hace $\Abar_t$ pequeño y el estado se reemplaza casi por la entrada ($t=3$: $h_3\approx x_3=-1$); un paso pequeño conserva el estado.
\end{solucion}

\begin{solucion}{eje:5-mit2}
\datos $y[n]=\frac12y[n-2]+x[n]$; $\Abar=\begin{pmatrix}0&1/2\\1&0\end{pmatrix}$, $\Bbar=(1,0)\T$, $C=(1,0)$.
\planteamiento (a) recurrencia con $x=\delta$; (b) potencias de $\Abar$ sobre $\Bbar$ y polinomio característico; (c) forma del núcleo de un SSM escalar.
\pasos
\begin{enumerate}[label=\arabic*.]
  \item (a) $y[0]=1$; $y[1]=\frac12y[-1]=0$; $y[2]=\frac12$; $y[3]=0$; $y[4]=\frac14$. Coincide con $K$.
  \item (b) $\Bbar=(1,0)\T$, $C\Bbar=1$; $\Abar\Bbar=(0,1)\T$, $C\cdot=0$; $\Abar^{2}\Bbar=(\frac12,0)\T$, $C\cdot=\frac12$; $\Abar^{3}\Bbar=(0,\frac12)\T$, $C\cdot=0$; $\Abar^{4}\Bbar=(\frac14,0)\T$, $C\cdot=\frac14$.
  \item $\det(\Abar-\lambda I)=\lambda^{2}-\frac12=0$, luego $\lambda=\pm 1/\sqrt2=\pm 0{,}7071$.
  \item (c) Un SSM escalar tiene $K_n=c\,\bar a^{\,n}\,\bar b$. $K_0=cb=1\ne 0$ obliga a $\bar a=0$ para que $K_1=cb\,\bar a=0$, pero entonces $K_2=0\ne\frac12$. Contradicción.
\end{enumerate}
\resultado \textbf{$C\Abar^{\,n}\Bbar=(1;\ 0;\ 0{,}5;\ 0;\ 0{,}25)$; polos $\pm 0{,}7071$; un SSM escalar no puede producir ceros alternos.} $\blacksquare$
\comprobacion $\Abar^{2}=\frac12I$, de modo que cada dos pasos el estado se multiplica por $\frac12$; los dos polos de signos opuestos se cancelan en los instantes impares.
\end{solucion}

\begin{solucion}{eje:5-barrido}
\datos $e_t=(a_t,b_t)$: $(0{,}5;\,1)$, $(0{,}5;\,0)$, $(0{,}5;\,0)$, $(0{,}5;\,2)$; $(a,b)\bullet(a',b')=(a'a,\ a'b+b')$.
\planteamiento agrupar en árbol, como haría un barrido paralelo, y usar la \cref{prop:scan}.
\pasos
\begin{enumerate}[label=\arabic*.]
  \item $e_1\bullet e_2=(0{,}25;\ 0{,}5\cdot 1+0)=(0{,}25;\ 0{,}5)$.
  \item $e_3\bullet e_4=(0{,}25;\ 0{,}5\cdot 0+2)=(0{,}25;\ 2)$.
  \item $(0{,}25;\ 0{,}5)\bullet(0{,}25;\ 2)=(0{,}0625;\ 0{,}25\cdot 0{,}5+2)=(0{,}0625;\ 2{,}125)$. Luego $h_4=2{,}125$.
  \item (b) $(1,0)\bullet(a,b)=(a\cdot 1,\ a\cdot 0+b)=(a,b)$ y $(a,b)\bullet(1,0)=(a,\ b+0)=(a,b)$.
  \item No conmutatividad: $(0{,}5;\,1)\bullet(1;\,2)=(0{,}5;\ 3)$, mientras que $(1;\,2)\bullet(0{,}5;\,1)=(0{,}5;\ 0{,}5\cdot 2+1)=(0{,}5;\ 2)$.
\end{enumerate}
\resultado \textbf{$h_4=2{,}125$, como en el \cref{ej:ssm}; $(1,0)$ es el neutro y $\bullet$ no es conmutativa, por lo que el barrido debe respetar el orden de los pasos.} $\blacksquare$
\comprobacion la primera componente, $0{,}0625=0{,}5^{4}$, es el factor por el que se multiplicaría un estado inicial no nulo $h_0$.
\end{solucion}

% ---------------------------------------------------------------------
\section{Tema 6}

\begin{solucion}{eje:6-parche}
\datos imagen $224\times 224$; parches de 16 o de 32.
\planteamiento \cref{prop:vit} y \cref{tab:vit}.
\pasos
\begin{enumerate}[label=\arabic*.]
  \item Parches: $(224/16)^{2}=196$ frente a $(224/32)^{2}=49$; cociente $4$.
  \item Coste de la atención: $\propto N_p^{2}$, se divide por $4^{2}=16$; capas lineales del Transformer: por 4.
  \item Dimensión de cada parche: $768$ frente a $32^{2}\cdot 3=3072$; la proyección $E$ crece 4 veces, pero se aplica a 4 veces menos parches.
  \item Se pierde resolución espacial: cada token resume una zona 4 veces mayor y los detalles pequeños quedan dentro de un único vector.
\end{enumerate}
\resultado \textbf{4 veces menos tokens y unas 16 veces menos coste de atención, a cambio de una representación más gruesa de la imagen.}
\comprobacion el coste de la proyección inicial no cambia: $196\cdot 768\cdot d=49\cdot 3072\cdot d=150\,528\,d$ multiplicaciones.
\end{solucion}

\begin{solucion}{eje:6-res}
\datos parche $16$; lados $224$ y $384$.
\planteamiento $N_p=(\text{lado}/16)^{2}$; lineales $\propto N_p$ y atención $\propto N_p^{2}$.
\pasos
\begin{enumerate}[label=\arabic*.]
  \item $224/16=14$, $N_p=196$; $384/16=24$, $N_p=576$.
  \item Cociente de tokens: $576/196=2{,}939$.
  \item Atención: $2{,}939^{2}=8{,}636$.
\end{enumerate}
\resultado \textbf{196 y 576 parches; las capas lineales crecen unas $2{,}939$ veces y la atención unas $8{,}636$.}
\comprobacion $(384/224)^{2}=1{,}714^{2}=2{,}939$ y $(384/224)^{4}=8{,}636$, como predice la \cref{prop:vit} con $\kappa=1{,}714$.
\end{solucion}

\begin{solucion}{eje:6-d2l}
\datos $P=16$, lados 96 y 192, token de clase, $d=512$, $\dff=2048$.
\planteamiento $n=(\text{lado}/P)^{2}+1$; lineales $n(8d^{2}+4d\dff)$; atención $4n^{2}d$.
\pasos
\begin{enumerate}[label=\arabic*.]
  \item $96/16=6$, $n=37$; $192/16=12$, $n=145$.
  \item Factor lineal: $8\cdot 262\,144+4\cdot 512\cdot 2048=2\,097\,152+4\,194\,304=6\,291\,456$ por token.
  \item $n=37$: lineales $232\,783\,872$; atención $4\cdot 1369\cdot 512=2\,803\,712$; total $235\,587\,584$.
  \item $n=145$: lineales $912\,261\,120$; atención $4\cdot 21\,025\cdot 512=43\,059\,200$; total $955\,320\,320$.
  \item Crecimiento: $955\,320\,320/235\,587\,584=4{,}055$. Fracción de atención: $1{,}190\,\%$ y $4{,}507\,\%$.
  \item Igualdad: $4n^{2}d=n\cdot 6\,291\,456\iff n=6\,291\,456/2048=3072$.
\end{enumerate}
\resultado \textbf{37 y 145 tokens; el coste por bloque crece $4{,}055$ veces y la atención pasa del $1{,}190\,\%$ al $4{,}507\,\%$; la atención dominaría a partir de $n=3072$ tokens.}
\comprobacion los tokens crecen $145/37=3{,}919$ y la atención $(145/37)^{2}=15{,}36$; el total queda entre ambos, cerca del primero, porque dominan las capas lineales.
\end{solucion}

\begin{solucion}{eje:6-whisper}
\datos $16$~kHz; fragmentos de $30$~s; desplazamiento de $10$~ms; convolución con paso~2; grabación de 10 minutos; códec de 50~Hz con 8 diccionarios.
\planteamiento contar muestras, fotogramas y vectores; para el códec, \cref{prop:rvq}.
\pasos
\begin{enumerate}[label=\arabic*.]
  \item (a) $16\,000\cdot 30=480\,000$ muestras.
  \item (b) $30/0{,}01=3000$ vectores de 80 canales; tras el paso~2, $1500$.
  \item (c) $600/30=20$ fragmentos; $20\cdot 1500=30\,000$ vectores.
  \item (d) Códec: $50\cdot 8=400$ tokens por segundo; en 600~s, $240\,000$ tokens, 8 veces más.
\end{enumerate}
\resultado \textbf{(a) $480\,000$ muestras; (b) 3000 y 1500 vectores; (c) $30\,000$ vectores; (d) $240\,000$ tokens con el códec.} El codificador de Whisper comprime 16 veces más que la sucesión de índices del códec, pero sus vectores son continuos y no sirven como objetivo de predicción del siguiente token.
\comprobacion cada vector del codificador cubre $600~\mathrm{s}/30\,000=0{,}02$~s $=20$~ms de audio, el doble del desplazamiento de $10$~ms del espectrograma, como corresponde al paso~2 de la convolución.
\end{solucion}

\begin{solucion}{eje:6-vitb}
\datos ViT-Base: $d=768$, $P=16$, $C_{\mathrm{in}}=3$, $N_p=196$, $8{,}6\cdot 10^{7}$ parámetros.
\planteamiento \cref{prop:vit}(i) y \cref{def:vit}.
\pasos
\begin{enumerate}[label=\arabic*.]
  \item $E$: $16^{2}\cdot 3\cdot 768=768\cdot 768=589\,824$.
  \item $E_{\mathrm{pos}}$: $(196+1)\cdot 768=151\,296$.
  \item Token de clase: $768$.
  \item Suma: $741\,888$; fracción $741\,888/8{,}6\cdot 10^{7}=0{,}8627\,\%$.
\end{enumerate}
\resultado \textbf{$589\,824+151\,296+768=741\,888$ parámetros, un $0{,}8627\,\%$ del modelo.} Casi todos los parámetros están en los bloques del Transformer, como en un modelo de lenguaje sin vocabulario grande.
\comprobacion cada bloque de ViT-Base tiene unos $4d^{2}+2d\cdot 3072=2\,359\,296+4\,718\,592=7\,077\,888$ parámetros de matriz; 12 bloques suman $8{,}49\cdot 10^{7}$, que con la entrada da $\approx 8{,}6\cdot 10^{7}$.
\end{solucion}

\begin{solucion}{eje:6-conv}
\datos imagen $X$ de lado $\ell$; filtros de $P\times P\times C_{\mathrm{in}}$; paso $S$.
\planteamiento escribir la salida de la convolución en la posición $(u,v)$ y compararla con la proyección del parche.
\pasos
\begin{enumerate}[label=\arabic*.]
  \item (a) Una convolución con filtro $F^{(c)}\in\R^{P\times P\times C_{\mathrm{in}}}$ y paso $S$ da en $(u,v)$ el valor $\sum_{a,b,\kappa}F^{(c)}_{a,b,\kappa}X_{uS+a,\,vS+b,\,\kappa}$.
  \item Con $S=P$, los índices $(uP+a,vP+b)$, $0\le a,b<P$, recorren exactamente el parche $(u,v)$, y los parches no se solapan.
  \item Si $x_p^{(u,v)}$ es el parche aplanado en el mismo orden que $F^{(c)}$, el valor es $x_p^{(u,v)}\cdot\mathrm{vec}(F^{(c)})$. Tomando como columna $c$ de $E$ el vector $\mathrm{vec}(F^{(c)})$, la salida en $(u,v)$ es la fila $x_p^{(u,v)}E$.
  \item (b) Las esquinas superiores de las ventanas son $0,S,2S,\dots,(\ell-P)$ en cada eje: $(\ell-P)/S+1$ posiciones, y en dos dimensiones, su cuadrado.
  \item Con $\ell=224$, $P=16$ y $S=8$: $(208/8+1)^{2}=27^{2}=729$.
\end{enumerate}
\resultado \textbf{(a) la incrustación por parches es una convolución con $d$ filtros $P\times P$ y paso $P$; (b) con paso $S$ hay $((\ell-P)/S+1)^{2}$ vectores: 729 en el ejemplo.} $\blacksquare$
\comprobacion con $S=P=16$, $(208/16+1)^{2}=14^{2}=196$, el número de parches de la \cref{tab:vit}.
\end{solucion}
